\section{Conclusion}


TLA+ model checking finds bugs in distributed protocols that no amount of testing will
catch. The bugs it finds are exactly the subtle ones: race conditions, deadlocks, and
safety violations that occur only under specific message orderings. Writing a TLA+
specification forces you to think precisely about your protocol's behavior in all
scenarios---not just the happy path.

For Lux Network, TLA+ model checking found 3 bugs in the Wave protocol and verified
the safety and conservation properties of the Teleport bridge across 2.4 million states.
These specifications now serve as the canonical descriptions of these protocols, and
every protocol change requires updating the specification and re-running TLC before
modifying the Go implementation.

\begin{thebibliography}{99}
\bibitem{lamport2002} Lamport, L. ``Specifying Systems: The TLA+ Language and Tools for Hardware and Software Engineers.'' Addison-Wesley, 2002.
\bibitem{lamport1994} Lamport, L. ``The Temporal Logic of Actions.'' ACM TOPLAS, 1994.
\bibitem{newcombe2015} Newcombe, C. et al. ``How Amazon Web Services Uses Formal Methods.'' Communications of the ACM, 2015.
\bibitem{yin2019} Yin, M. et al. ``HotStuff: BFT Consensus with Linearity and Responsiveness.'' PODC, 2019.
\bibitem{castro1999} Castro, M. and Liskov, B. ``Practical Byzantine Fault Tolerance.'' OSDI, 1999.
\bibitem{tlc} Yu, Y., Manolios, P., and Lamport, L. ``Model Checking TLA+ Specifications.'' CHARME, 1999.
\end{thebibliography}

